\documentclass[10pt]{article}
% ==================  PACKAGES ==================
\usepackage{graphicx}
\usepackage{amssymb}
\usepackage{amsmath}
\usepackage{amsthm}
\usepackage{mathtools}
\usepackage{verbatim}
\usepackage{url}



% ==================  SETTINGS  ==================
\setlength{\textwidth}{5.325in}
\setlength{\textheight}{8.25in}
\setlength{\topmargin}{-0.4cm}
\setlength{\evensidemargin}{1.6cm}
\setlength{\oddsidemargin}{1.6cm}
\UseRawInputEncoding % Forces compatiblitity with new latex UTF encoding



% ==================  SHANE MATH DEFINITIONS  ==================
% ---------------------- bb fonts ---------------------- 
\newcommand{\rr}{\mathbb{R}}
\newcommand{\qq}{\mathbb{Q}}
\newcommand{\cc}{\mathbb{C}}
\newcommand{\zz}{\mathbb{Z}}
\newcommand{\nn}{\mathbb{N}}
\newcommand{\dd}{\mathbb{D}}
\newcommand{\pp}{\mathbb{P}}
% ----------------------- cal fonts ----------------------
\newcommand{\cP}{\mathcal{P}}
\newcommand{\cT}{\mathcal{T}}
\newcommand{\cX}{\mathcal{X}}
\newcommand{\cQ}{\mathcal{Q}}
\newcommand{\cF}{\mathcal{F}}
\newcommand{\bigo}{\mathcal{O}}
% ---------------------- interval arithmetic ---------------------- 
\newcommand{\intrr}{\mathbb{IR}}
\newcommand{\intcc}{\mathbb{IC}}
\newcommand{\intval}[1]{\tilde{#1}}
\DeclareMathOperator{\rad}{rad}
% ---------------------- variable flair ---------------------- 
\newcommand{\overbar}[1]{\mkern 1.5mu\overline{\mkern-1.5mu#1\mkern-1.5mu}\mkern 1.5mu}
% ---------------------- operator names ---------------------- 
%\DeclareMathOperator{\ker}{ker}
\DeclareMathOperator{\coker}{coker}
\DeclareMathOperator{\image}{image}
\DeclareMathOperator{\sign}{sign}
\DeclareMathOperator{\inspan}{span}
% ---------------------- delimeters ---------------------- 
\newcommand{\setof}[1]{\left\{#1\right\}}
\newcommand{\pardiff}[2]{\frac{\partial #1}{\partial #2}}
\newcommand{\abs}[1]{\left| #1 \right|}
\DeclarePairedDelimiter{\norm}{\lVert}{\rVert}
\DeclarePairedDelimiter{\dotp}{\langle}{\rangle}
%\newcommand{\dotp}[1]{\langle #1 \rangle}
\newcommand{\paren}[1]{\left( #1 \right)}
% ---------------------- misc ---------------------- 
\renewcommand{\Re}{\text{Re}}
\renewcommand{\Im}{\text{Im}}
\renewcommand{\emptyset}{\O}
\newcommand{\ptlt}{\hskip 5pt < \hskip-7pt \cdot} % pointwise less than
\newcommand{\ptgt}{\hskip 5pt \cdot \hskip-7pt >} % pointwise greater than
\newcommand{\bydef}{\,\stackrel{\mbox{\tiny\textnormal{\raisebox{0ex}[0ex][0ex]{def}}}}{=}\,}



% ==================  THEOREM ENVIRONMENTS ==================
\newtheorem{theorem}{Theorem}%[section]
\newtheorem{definition}{Definition}
\newtheorem{corrolary}[theorem]{Corollary}
\newtheorem{lemma}[theorem]{Lemma}
\newtheorem{proposition}[theorem]{Proposition}
\newtheorem{remark}{Remark}
\newtheorem{example}{Example}
%\numberwithin{equation}{section}
% ---------------------- theorem list environment ---------------------- 
\newcounter{listacnt}\renewcommand{\thelistacnt}{\alph{listacnt}}
\newenvironment{theoremlist}{\begin{list}{\textnormal{(\thelistacnt)}}%
		{\settowidth{\labelwidth}{(b)}%
			\setlength{\topsep}{1.5pt plus 0.5pt minus 0.5pt}%
			\setlength{\itemsep}{1pt plus 0.2pt minus 0.2pt}%
			\setlength{\parsep}{0.1pt plus 0.1pt minus 0.1pt}%
			\usecounter{listacnt}}}{\end{list}}



% ==================  START OF DOCUMENT ==================
\begin{document}
\title{CRTBP Collision manifold integration notes}%
\author{Shane Kepley}
\date{}
\maketitle


\section{Choosing the best coordinates}

We fix our parameter, $\mu$, to be the mass of the small primary with is located at the position $(\mu-1, 0$). So we have $0 < \mu \leq \frac{1}{2}$ and similarly, the large primary has mass $1 - \mu$ and its position in configuration space is $(\mu, 0)$.
\subsection{The 3 vector fields}
Once and for all we write down the formulas for all 3 vector fields. Here $\mu_2 = \mu$ and  $\mu_1 = 1- \mu_2$ in any case they appear. We let $f_0, f_1, f_2$ denote the CRTBP vector field, the regularized field at the large primary, and the regularized field at the small primary respectively.
\subsubsection{Standard}
The Szebehely version of the CRTBP has mass parameter $0 < \mu \leq \frac{1}{2}$. The large primary is located at $P_1 = (\mu, 0)$ having mass, $1 - \mu$ and the small primary is located at $P_2 = (\mu - 1, 0)$ having mass, $\mu$. The vector field prior to automatic differentiation is given by the equations
\begin{align*}
	\dot x & = p \\
	\dot p & = 2q + \left(1 - \mu\right)\left(x - \mu\right) + \mu\left(x + 1 - \mu\right) - \frac{\paren{1 - \mu} \paren{x - \mu}}{d_1^3} - \frac{\mu \paren{x + 1 - \mu}}{d_2^3} \\
	\dot y &= q \\
	\dot q & = - 2p + y - \frac{\paren{1-\mu}y}{d_1^3} - \frac{\mu y}{d_2^3}
\end{align*} 
where $d_1, d_2$ are defined by the formulas
\[
d_1 = \sqrt{\paren{x - \mu}^2 + y^2} \qquad d_2 = \sqrt{\paren{x + 1 - \mu}^2 + y^2}
\]
Below, we will often write $\mu_1 = 1 - \mu$ and $\mu_2 = \mu$. 

The (extended) standard CRTBP has coordinates are denoted by $(x_0,p_0,y_0,q_0,r_0,s_0))$. The vector field is denoted by, $f_0 : \rr^6 \to \rr^6$, is obtained after automatic differentiation by setting $r_0 = d_1^{-1}$ and $s_0 = d_2^{-1}$. The explicit formulas are now given by the polynomial formulas
\begin{align*} 
\dot x_0 & = p_0 \\ 
\dot p_0 & = 2 q_0 + \mu _ { 1 } \left( x_0 - \mu _ { 2 } \right) + \mu _ { 2 } \left( x_0 + \mu _ { 1 } \right) - \mu _ { 1 } \left( x_0 - \mu _ { 2 } \right) r_0^3 - \mu _ { 2 } \left( x_0 + \mu _ { 1 } \right) s_0^3 \\ 
\dot y_0 & = q_0 \\ 
\dot q_0 & = - 2 p_0 + \mu _ { 1 } y_0 + \mu _ { 2 } y_0 - \mu _ { 1 } y_0 r_0^3 - \mu _ { 2 } y_0 s_0^3 \\
\dot{r_0} & = - r_0^3 \left( \left( x_0 - \mu_2 \right) p_0 + y_0 q_0 \right) \\ 
\dot{s_0}  & = - s_0^3 \left( \left( x_0 + \mu_1\right) p_0 + y_0 q_0 \right) 
\end{align*}
where $r_0$ and $s_0$ are the automatic differentiation coordinates given by 
\[
r_0 = \frac{1}{\sqrt{(x_0-\mu_2)^2 + y_0^2}} \qquad s_0 = \frac{1}{\sqrt{(x_0+\mu_1)^2 + y_0^2}}.
\]
After factoring, it is often convenient to write the formulas for $\dot p_0, \dot q_0$ as 
\begin{align*}
	\dot p_0 & = 2 q_0 + \mu_1 \paren{x - \mu_2} \paren{1 - r_1^3} + \mu_2 \paren{x + \mu_1} \paren{1 - r_2^3} \\
	\dot q_0 & = -2 p_0 + \mu_1 y \paren{1 - r_1^3} + \mu_2 y \paren{1 - r_2^3}
\end{align*}

\subsubsection{Large primary regularization}
The vector field $f_1: \rr^5 \to \rr^5$ is the vector field desingularized at the large primary located at coordinates $(\mu,0)$ in configuration space. It provides the derivatives given by the formulas
\begin{align*}
\dot x_1 & = p_1 \\
\dot p_1 & = { 8 \left( x_1^2 + y_1^2 \right) q_1 + 12 x_1 \left( x_1^2 + y_1^2 \right)^2 + 16 \mu x_1^3 + 4 ( \mu - C ) x_1 + 8 \mu x_1 \left( x_1^2 - 3y_1^2 + 1 \right) r_1^3 } \\
\dot y_1 & = q_1 \\
\dot q_1 & =  { - 8 \left( x_1^2 + y_1^2 \right) p_1 + 12 y_1 \left( x_1^2 + y_1^2 \right)^2 - 16 \mu y_1^3 + 4 ( \mu - C ) y_1 + 8 \mu y_1 \left( - y_1^2 + 3 x_1^2+ 1 \right) r_1^3 } \\
 \dot r_1 & =   { - 2 r_1^3 \left( \left( x_1^2 + y_1^2 \right) ( x_1 p_1 + y_1 q_1 ) + ( x_1 p_1 - y_1 q_1 ) \right) }
\end{align*}
where $C = H_0(x_0,p_0,y_0,q_0)$ is the standard energy and $r_1$ is given by 
\[
r_1 = \frac{1}{\sqrt{(x_1^2 + y_1^2)^2 + 1 + 2(x_1^2 - y_1^2)}}.
\]


\subsubsection{Small primary regularization}
The vector field $f_2: \rr^5 \to \rr^5$ is the vector field desingularized at the small primary located at coordinates $(\mu-1,0)$. It provides the derivatives given by the formulas
\begin{align*} 
\dot x_2 & = p_2 \\ 
\dot p_2 & = 8 \left( x_2^2 + y_2^2 \right) q_2 + 12 x_2 \left( x_2^2 + y_2^2  \right)^2 + 16(\mu-1) x_2^3 + 4 ( 1-\mu - C) x_2 + 8 ( 1 - \mu ) x_2 \left( - x_2^2+ 3y_2^2 + 1 \right) r_2^3 \\
\dot y_2 & = q_2 \\ 
\dot q_2 & = - 8 \left( x_2^2 + y_2^2 \right) p_2 + 12 y_2 \left( x_2^2 + y_2^2 \right)^2 - 16 (\mu-1 ) y_2^3 + 4 ( 1-\mu-C ) y_2 + 8 ( 1 - \mu )y_2 \left( y_2^2 - 3 x_2^2 + 1 \right)r_2^3\\
\dot{r_2}  & =  - 2 r_2^3 \left( \left( x_2^2 + y_2^2\right) ( x_2 p_2 + y_2 q_2 ) - ( x_2 p_2 - y_2 q_2 ) \right) 
\end{align*}
where $C = H_0(x_0,p_0,y_0,q_0)$ is the standard energy and $r_2$ is given by 
\[
r_2 = \frac{1}{\sqrt{(x_2^2 + y_2^2)^2 + 1 + 2(y_2^2 - x_2^2)}}.
\]


\subsection{Changing coordinates}

\subsection{Mappings between coordinates}
The basis for these coordinate transformations are the mappings $z \mapsto w^2 + \mu$ and $z \mapsto w^2 + \mu - 1$ which define the desingularized coordinates about the large and small primaries respectively. 
\subsubsection{standard to large primary}
For the regularized vector field, $f_1$ (desingularized at the large primary) we have the coordinates $(x_1,p_1,y_1,q_1)$ which correspond to the standard CRTBP coordinates $(x_0,p_0,y_0,q_0)$ by the formulas: 
\begin{align*}
\label{eq:f1_f0}
	x_0 & = x_1^2 - y_1^2 + \mu \\
	p_0 & = \frac{x_1p_1 - y_1q_1}{2(x_1^2 + y_1^2)} \\
	y_0 & = 2x_1y_1 \\
	q_0 & = \frac{y_1p_1 + x_1q_1}{2(x_1^2 + y_1^2)}
\end{align*}
Note that $p_0,q_0$ are the $t$-time derivatives of $x_0,y_0$ while $p_1,y_1$ are the $\tau$-time derivatives of $x_1,y_1$ and the time rescaling is related by 
\[
\frac{dt}{d \tau} = 4w \overbar{w} \qquad   \frac{d \tau }{dt} = \frac{1}{4 w \overbar{w}}
\]
where $w \overbar{w} = x_1^2 + y_1^2$.
Then we have for example from $z = x_0 + iy_0 = w^2 + \mu$, the formula 
\begin{align*}
\frac{dz}{dt} & = 2 w \frac{dw}{d \tau}\frac{d\tau}{dt} \\
& = 2(x_1 + iy_1)(p_1 + iq_1)(\frac{1}{4 w \overbar{w}}) \\
& = \frac{x_1p_1 - y_1q_1}{2(x_1^2 + y_1^2)} + i \frac{x_1q_1 + y_1p_1}{2(x_1^2 + y_1^2)}
\end{align*}
and the values of $p_0,q_0$ follow from noting that $\dot z = \dot x_0 + i\dot y_0$.

These formulas can be solved in the other direction to obtain the regularization transformation
\begin{align*}
	x_1 & =  \Re\left(\sqrt{x_0 + iy_0 - \mu}\right) \\
	p_1 & = 2(p_0x_1 + q_0y_1) \\
	y_1 & = \Im\left(\sqrt{x_0 + iy_0 - \mu}\right) \\
	q_1 & = 2(q_0x_1 - p_0y_1)
\end{align*}
Again, this computation follows from the computation
\[
\frac{d}{d \tau}(w^2) = 2w \frac{dw}{d \tau} = \frac{d}{d \tau} (z - \mu) = \frac{dz}{dt} \frac{dt}{d \tau}.
\]
So we obtain 
\begin{align*}
\frac{dw}{d \tau} & = \frac{4 \dot z w \overbar w}{2w} \\
& =2 \dot z \overbar w \\
& = 2(p_0 x_1 + q_0 y_1) + 2i(q_0 x_1 - p_0 y_1). 
\end{align*}
The potential in the new coordinates is 
\[
\mathcal{U}_1(x_1,y_1) =  2 ( x_1^2 + y_1^2 )^3 + 4\mu ( x_1^4 - y_1^4 ) + 2( \mu - C ) ( x_1^2 + y_1^2 ) + 4 ( 1 - \mu ) + \frac { 4 \mu ( x_1^2 + y_1^2 ) } { \sqrt { ( x_1^2 + y_1^2 )^2 + 1 + 2 ( x_1^2 - y_1^2 ) } }.
\]
which gives rise to the the Jacobi constant in the new coordinates which is
\[
H_1(x_1,p_1,y_1,q_1) = 2\mathcal{U}_1(x_1,y_1) - p_1^2 - q_1^2.
\]

\subsubsection{standard to small primary}
Working similar to the previous computations, we have formulas for the small body regularization. 
\begin{align*}
x_0 & = x_2^2 - y_2^2 + \mu-1\\
p_0 & = \frac{x_2p_2 - y_2q_2}{2(x_2^2 + y_2^2)} \\
y_0 & = 2x_2y_2 \\
q_0 & = \frac{x_2q_2 + y_2p_2}{2(x_2^2 + y_2^2)} \\
\end{align*}
These formulas can be solved in the other direction to obtain the regularization transformation
\begin{align*}
x_2 & = \Re\left(\sqrt{x_0 + iy_0 - (\mu-1)}\right) \\
p_2 & =  2 \left(p_0x_2 + q_0 y_2 \right)\\
y_2 & = \Im\left(\sqrt{x_0 + iy_0 - (\mu-1)}\right) \\
q_2 & = 2 \left(q_0x_2 - p_0 y_2\right)
\end{align*}
The potential in the new coordinates is 
\[
\mathcal{U}_2(x_2,y_2) =  2 (x_2^2 + y_2^2 )^3 + 4(\mu-1) (x_2^4 - y_2^4) + 2(1-\mu - C) (x_2^2 + y_2^2) + 4\mu + \frac { 4 (1-\mu) ( x_2^2 + y_2^2 ) } { \sqrt { (x_2^2 + y_2^2 )^2 + 1 + 2 ( x_2^2 - y_2^2 ) } }.
\]
which gives rise to the the Jacobi constant in the new coordinates which is
\[
H_2(x_2,p_2,y_2,q_2) = 2\mathcal{U}_2(x_2,y_2) - p_2^2 - q_2^2.
\]

We summarize and denote these changes of coordinates in a vectorized form for further reference. Begin by defining three Euclidean spaces, $X_i \cong \rr^4$ for $i = 0,1,2$. Of course these spaces are isomorphic, however, we will write $v \in X_i$ when referring to points which are being integrated by the corresponding vector field, $f_i$. In other words, $X_i$ will be used to denote the phase space for the dynamics generated by $f_i$. We define maps between these spaces, $\mathcal{F}_i : X_0 \to X_i$ for $i = 1,2$ given by 
\[
(x_0,p_0,y_0,q_0) \overset{\mathcal{F}_i}{ \longmapsto} (x_i,p_i,y_i,q_i)
\]
and similarly, the inverse maps, $F_i^{-1}: X_i \to X_0$ are defined 
\[
(x_i,p_i,y_i,q_i)  \overset{\mathcal{F}^{-1}_i}{ \longmapsto}   (x_0,p_0,y_0,q_0).
\]

\subsubsection{small to large primary and back}
These maps naturally compose to produce a map between regularized coordinates. For example, changing from $f_2$ to $f_1$ coordinates corresponds to the map
\[
(x_1 ,p_1 ,y_1 ,q_1) = \mathcal{F}_1 \circ \mathcal{F}^{-1}_2 \left(x_2,p_2,y_2,q_2\right).
\]
Using the above formulas for $\cF_1, \cF_2$, we can write the formulas for $\mathcal{F}_1 \circ \mathcal{F}^{-1}_2$ which is given by
\begin{align*}
x_1 & =  \Re \left(\sqrt{x_2^2 - y_2^2 - 1 + i(2x_2y_2) } \right) \\
p_1 & = \frac{x_1(x_2p_2 - y_2q_2) + y_1(x_2q_2 + y_2p_2)}{x_2^2 + y_2^2} \\
y_1 & = \Im \left(\sqrt{x_2^2 - y_2^2 - 1 + i(2x_2y_2) } \right) \\
q_1 & = \frac{x_1(x_2q_2 + y_2p_2) - y_1(x_2p_2 - y_2q_2)}{x_2^2 + y_2^2}
\end{align*}
Likewise, we obtain formulas for the map from $f_1$ coordinates to $f_2$ coordinates given by $\mathcal{F}_2 \circ \mathcal{F}^{-1}_1$ given by 
\begin{align*}
x_2 & =  \Re \left( \sqrt{x_1^2 - y_1^2 + 1 + i(2x_1y_1) } \right) \\
p_2 & = \frac{x_2(x_1p_1 - y_1q_1) + y_2(x_1q_1 + y_1p_1)}{x_1^2 + y_1^2} \\
y_2 & = \Im \left( \sqrt{x_1^2 - y_1^2 + 1 + i(2x_1y_1) } \right) \\
q_2 & = \frac{x_2(x_1q_1 + y_1p_1) - y_2(x_1p_1 - y_1q_1)}{x_1^2 + y_1^2}
\end{align*}


\subsubsection{Normalizing time}
In order to keep a global track of time for orbits with respect to their position in $X_0$, we will normalize the time along orbits in regularized coordinates. We let $t \in \rr$ always denote this time and use $\tau_1, \tau_2$ respectively, to keep time along orbits when integrating in $X_1,X_2$ respectively. This normalization must satisfy the following property. Suppose $v \in X_0$ is an arbitrary initial condition. Let $\gamma_0$ denote the trajectory through $v$ satisfying $\gamma_0(0) = v$. Now, let $\gamma_i$ denote the trajectory segment for $f_i$ satisfying $\gamma_i(0) = \mathcal{F}_i(v)$. Then, the normalized time along $\gamma_i$ is the mapping, $t : \rr \to \rr$ satisfying
\[
\gamma_0(t(\tau)) = \mathcal{F}_i^{-1} \left( \gamma_i(\tau)\right) \qquad \text{for all} \ \tau \in \rr. 
\]
Recalling the coordinate changes defined $\mathcal{F}_i$ above, we have $\frac{dt}{d \tau_i} = 4w_i \overbar{w}_i$. Noting that $w_i \overbar{w}_i = \abs{w_i}^2 = x_i^2 + y_i^2$, we can directly solve this differential equation to obtain the formula
\[
t(\tau) = 4\int_0^{\tau} x_i(\tau_i)^2 + y_i(\tau_i)^2 \ d \tau_i .
\]
Similarly, we have $\frac{d \tau_i}{dt} = \frac{1}{4 w_i \overbar{w}_i}$


\subsection{When to change coordinates}
We are interested in propagating material surfaces by dynamically switching to the coordinates which have the best properties.  In our case, we would like to obtain the largest surface area per time-step for the Taylor integrator and this amounts to remaining bounded away from the singular points of each vector field. An analogous property which is easy to measure is the (squared) {\em speed} which we define pointwise in any coordinates as $S_i = p_i^2 + q_i^2 $ for $ i = 0,1,2$. 

A reasonable strategy then is to integrate using whichever of the 3 vector fields has the minimum value for this speed along ``most'' of the boundary chart. To determine this, we compare $S_1 = p_1^2 + q_1^2$ in the (large primary) regularized coordinates, to the corresponding speed in the standard CRTBP coordinates, which is given explicitly by 
\begin{align*}
S_0 = p_0^2 + q_0^2  & = \frac{(x_1p_1-y_1q_1)^2 + (y_1p_1 + x_1q_1)^2}{4(x_1^2 + y_1^2)^2} \\
& = \frac{(x_1p_1)^2 + (y_1q_1)^2 + (y_1p_1)^2 + (x_1q_1)^2}{4(x_1^2 + y_1^2)^2} \\
& = \frac{(x_1^2 + y_1^2)(p_1^2 + q_1^2)}{4(x_1^2 + y_1^2)^2} \\ 
& = \frac{(p_1^2 + q_1^2)}{4(x_1^2 + y_1^2)}
\end{align*}
Evidently, we have $S_0 = S_1$ if and only if $x_1^2 + y_1^2 = \frac{1}{4}$ which is a circle of radius $\frac{1}{2}$ centered at the origin in the $(x_01y_1)$ coordinates. We call this circle the {\em ideal} domain for $f_1$. Hence, for any points inside this circle we have $S_1 < S_0$ and we integrate using the regularized vector field. Otherwise, we switch to the standard coordinates before integrating and the formula $x_1^2 + y_1^2 = \frac{1}{4}$ provides the criteria to be checked when determining whether or not to change coordinates. 

Similarly, if we are in $f_0$ coordinates, we would like a formula for determining when to change to $f_1$ coordinates. For this we compute the image of the ideal circle for $f_1$ using the formulas for $x_0,y_0$ above. In anticipation that this domain is a circle centered at $(\mu,0)$ in the $(x_0,y_0)$ plane, we compute
\begin{align*}
	(x_0 - \mu)^2 + y_0^2 & = (x_1^2 - y_1^2)^2 + (2x_1y_1)^2 \\
	& = (x_1^2 + y_1^2)^2 \\
	& = \frac{1}{16}
\end{align*}
where in the last step we used the assumption that $x_1^2 + y_1^2 = \frac{1}{4}$. Thus, the ideal domain for which $S_0 < S_1$ is the interior of the disc bounded by the circle of radius $\frac{1}{4}$ centered at the large primary in the $(x_0,y_0)$ plane. We use the formula, $(x_0-\mu)^2 + y_0^2 = \frac{1}{16}$ jas the computable criteria for switching from $f_0$ coordinates into $f_1$ coordinates. 


A similar computation shows that for points written in $f_2$ coordinates (near the small primary) there is a disc of radius $\frac{1}{2}$ given by the formula $x_2^2 + y_2^2 = \frac{1}{4}$ on which $S_2 < S_0$. This is the ideal domain for $f_2$ and we flow using this vector field inside of this disc, and otherwise, we switch to the $f_0$ coordinates. The corresponding condition for switching from $f_0$ into $f_2$ coordinates is given by the  boundary circle $(x_0 - \mu + 1)^2 + y_0^2 = \frac{1}{16}$. 

\subsection{Integration details}
Consider an arbitrary arc of initial data $\gamma: [-1,1] \to \cc^5$ with respect to, $f_i$, one of the regularized vector fields. We assume it is analytic with Taylor expansion given by 
\[
\gamma(\sigma) = \sum_{n = 0}^{\infty} a_n \sigma^n \qquad a_n \in \cc^5, \quad \sigma \in [-1,1].
\]
We also denote the equivalent coordinate-wise representation as follows. 
\[
\gamma(\sigma) = 
\begin{pmatrix}
x_i(\sigma) \\
p_i(\sigma) \\
y_i(\sigma) \\
q_i(\sigma) \\
r_i(\sigma)
\end{pmatrix}
=
\sum_{n = 0}^{\infty} 
\begin{pmatrix}
a^{(1)}_n \\
a^{(2)}_n \\
a^{(3)}_n \\
a^{(4)}_n\\
a^{(5)}_n
\end{pmatrix}
\sigma^n 
\]
That is, $x_i,p_i,y_i,q_i,r_i$ are considered as one variable power series. In this work, these parameterize portions of collision manifolds for one of the primary bodies, or invariant manifolds attached to one of the equilibria. For computational purposes, these are represented in the computer as truncated coefficient sequence which defines a polynomial, and rigorous error bounds for the analytic tail i.e.~an $\ell_1$ error bound on the truncation. We denote this by the representation $\gamma(\sigma) = (\overbar{\gamma}(\sigma),r)$ where $\overbar{\gamma}: \rr \to \cc^5$ is a polynomial of order $N \geq 1$ and $ r > 0$ satisfies the bound
\[
\max \setof{
\sum_{N+1}^{\infty} \abs{a^{(1)}_n}, \sum_{N+1}^{\infty} \abs{a^{(2)}_n}, \sum_{N+1}^{\infty} \abs{a^{(3)}_n}, \sum_{N+1}^{\infty} \abs{a^{(4)}_n}, \sum_{N+1}^{\infty} \abs{a^{(5)}_n}} < r < \infty. 
\]


\subsection{configuration space} 
We start with the configuration variables, $x,y$. We start by assuming $u(\sigma) = \overbar{u}(\sigma) + h_u(\sigma)$ and $v(\sigma) = \overbar{v}(\sigma) + h_v(\sigma)$ where $\overbar{u}, \overbar{v}$ are the polynomial parts and $h_u,h_v \in C^\omega(\dd)$ satisfy 
\[
\norm{\cT(h_u)}_{\ell_1} < r \qquad \norm{\cT (h_v)}_{\ell_1} < r. 
\]
From the formula for $x$ we have $x(\sigma) = u(\sigma)^2 - v(\sigma)^2 + \mu_2$ so the polynomial part is obtained as the vector $\overbar{x} = (x_0,\dots,x_N)$ as the first $N+1$ terms of  we have the bound for the Taylor transform of $x$ given by 
\begin{align*}
	\cT(x) & = \left( (u_0,\dots,u_N) + \cT(h_u) \right)^2 + \left( (u_0,\dots,u_N) + \cT(h_u) \right)^2 
\end{align*} 

\section{Regularized Jacobi Constant}
We also compute for later reference, the regularized energy function. Given the original energy and potential functions
\[
E(x,p,y,q) = \frac{1}{2}(p^2 +q^2) - \Omega(x,y) \qquad \Omega(x,y) = \frac{1}{2}(x^2 + y^2) + \frac{\mu_1}{r_0} + \frac{\mu_2}{s_0}.
\]


In the regularized (at $\mu_1$) coordinates we have $\mu = \mu_2$ so the potential is 
\[
\mathcal{U}(u,v) =  2 \left( u^2 + v^2 \right)^3 + 4 \mu \left( u^4 - v^4 \right) + 2( \mu - C ) \left( u^2 + v^2 \right) + 4 ( 1 - \mu ) + \frac { 4 \mu \left( u^2 + v^2 \right) } { \sqrt { \left( u^2 + v^2 \right)^2 + 1 + 2 \left( u^2 - v^2 \right) } }.
\]
We set $k_1 = u^2 + v^2$ and $k_2 = u^2 - v^2$ and rewrite this potential as
\[
\mathcal{U}(u,v) =  2 k_1^ { 3 } + 4 \mu k_1k_2+ 2( \mu - C )k_1 + 4 ( 1 - \mu ) + \frac { 4 \mu k_1} { \sqrt { k_1^2 + 1 + 2 k_2 } }.
\]
Now, we have the formulas for the partial derivatives
\begin{align*}
	\pardiff{\mathcal{U}}{u} & = 12 {\left(u^{2} + v^{2}\right)}^{2} u +  16 \mu u^3 + 4  {\left(\mu - C \right)} u + \frac{8 \mu u (u^2-3v^2+1)}{{\left({\left(u^{2} + v^{2}\right)}^{2} + 2(u^2-v^2) + 1\right)}^{\frac{3}{2}}} \\
	& = 12k_1^2 u +  16 \mu u^3 + 4  {\left(\mu - C \right)} u + \frac{8 \mu u (u^2-3v^2+1)}{(k_1^2 + 2k_2 + 1)^{\frac{3}{2}}} 
\end{align*}
and 
\begin{align*}
 \pardiff{\mathcal{U}}{u}  & = 12{\left(u^{2} + v^{2}\right)}^2 v - 16\mu v^3 + 4 {\left(\mu -C\right)} v  + 
 \frac{8 \mu v(3u^2 - v^2 +1)}{\left((u^2+v^2)^2 + 2(u^2-v^2) + 1 \right)^{\frac{3}{2}}} \\
 & = 12 k_1^2 v - 16 \mu v^3 + 4(\mu - C)v + \frac{8 \mu v(3u^2 - v^2 +1)}{(k_1^2 + 2k_2 + 1)^{\frac{3}{2}}} 
\end{align*}
The Jacobi constant is now given by the formula
\[
E(u,p,v,q) = p^2 + q^2 - 2\mathcal{U}(u,v) 
\]
We can evaluate $\frac{dE}{dt}$ along a trajectory using the formulas from the regularize vector field
\begin{eqnarray*}
\dot{u} = p \\
\dot{p} = 8k_1q + \pardiff{\mathcal{U}}{u} \\ 
\dot{v} = q \\
\dot{q} = -8k_1 p + \pardiff{\mathcal{U}}{v}.
\end{eqnarray*}
Now, we introduce the variables $D = \sqrt{k_1^2 + 2k_2 + 1}, h_1 = pu + qv$ and $h_2 = pu-qv$ and simplifying with $k_i,h_i,D$ we have 
\subsection{Formulas for kinetic energy}   
Lets compute $p \dot p$ and $q \dot q$. 
\begin{align*}
p \dot p & = 8k_1qp + 12k_1^2 up + 16 \mu u^3 p + 4(\mu - C) up + \frac{8 \mu u p}{D^3} \left(u^2 - 3v^2 + 1\right) \\
q \dot q & = -8k_1 pq + 12k_1^2 vq - 16 \mu v^3 q + 4(\mu-C) vq + \frac{8 \mu vq}{D^3} \left(3u^2 - v^2 + 1\right))
\end{align*}
and adding them together we get
\[
p \dot p + q \dot q= 12k_1^2 h_1 + 16 \mu (u^3 p - v^3 q) + 4(\mu - C) h_1 + \frac{8 \mu }{D^3} \left(u^3p - 3uv^2p + 3u^2 vq - v^3q + h_1\right).
\]
On the other hand, we have $\dot{k_1} = 2h_1$ and $\dot{k_2} = 2h_2$ and $\dot D = \frac{k_1 \dot k_1 + \dot k_2}{D}$ so we can compute $\frac{d \mathcal{U}}{dt}$ as
\begin{align*}
\dot {\mathcal{U}} & = 6k_1^2 \dot k_1 + 4 \mu(k_1 \dot k_2 + \dot k_1 + k_2) + 2(\mu - C) \dot k_1 + 4 \mu \left(\frac{\dot k_1}{D} - \frac{ k_1 \dot D}{D^2}  \right) \\ 
& = 12 k_1^2 h_1 + 8 \mu (k_1h_2 + h_1k_2) + 4(\mu - C) h_1 + 8 \mu (\frac{h_1}{ D} + \frac{k_1^2 h_1 + k_1 h_2}{D^3}).
\end{align*}
Note that the first and third terms cancel immediately. Also, we compute $k_1h_2 + h_1k_2 = 2(pu^3 - qv^3)$ so the third terms also cancel. To check that the final terms cancel, we need to verify that
\[
 u^3p - 3uv^2p + 3u^2 vq - v^3q + h_1 = h_1D^2 + k_1^2h_1 + k_1h_2.
\]

\section{Verifying connections with BVP}

\subsection{The Newton map without regularization}
We first consider the case of an approximate connecting orbit which is homoclinic to an equilibrium. In this case, despite our mining procedure utilizing the regularization technique, the verification of a connecting orbit does not require any regularization. The appropriate BVP fomrulation is as follows. 

Suppose $\hat x$ is a saddle-focus equilibrium for the CRTBP vector field, $f: \rr^4 \to \rr^4$ with a two-dimensional stable and unstable manifolds. We assume we have parameterized these manifolds. Specifically, we assume  $P : \cc^2 \to \rr^4$ and $Q: \cc^2 \to \rr^4$ satisfying $P(\dd^2) \subset W^u(\hat x)$ and $P(\dd^2) \subset W^s(\hat x)$ and $P, Q$ have been explicitly computed as high order Taylor expansions. Moreover, the Taylor coefficients for $P,Q$ have the conjugate symmetry property as defined in \ref{ } implying that a {\em real} local unstable manifold may be obtained explicitly as 
\[
W^u_{\text{loc}}(\hat x) = \setof{P\left(re^{i \theta}, r e^{-i\theta} \right) : 0 \leq r < 1, 0 \leq \theta < 2 \pi}.
\]
Moreover, we the boundary of the local unstable manifold is a $1$-dimensional manifold with an explicit parameterization given by the formula, 
\[
P_{\partial}(\theta) = P(e^{i \theta}, e^{-i \theta}) \quad \theta \in [0, 2\pi).
\]
Similarly, we have parameterizations for the local stable manifold and its boundary given by 
\begin{align*}
W^s_{\text{loc}}(\hat x) &= \setof{Q\left(re^{i \theta}, r e^{-i\theta} \right) : 0 \leq r < 1, 0 \leq \theta < 2 \pi} \\
Q_{\partial}(\theta)& = Q(e^{i \theta}, e^{-i \theta}) \quad \theta \in [0, 2\pi).
\end{align*}

Now, the Newton map for the multiple shooting scheme is defined to be $g: \rr^{4n-1} \to \rr^{4n}$ defined by 
\[
g(\alpha, x_1, \dotsc, x_{n-1}, \beta, \tau) = \left(g_1, \dotsc, g_n\right)
\]
where, $\setof{\alpha, \beta, \tau} \subset \rr$, $\setof{x_j : 1\leq j \leq n-1} \subset \rr^4$, and the coordinate functions are defined by 
\begin{align*}
g_1(\alpha, x_1) & = \Phi\left(t_1, P_\partial(\alpha)\right) - x_1 \\
g_j(x_{j-1}, x_j) & = \Phi \left(t_j, x_{j-1}\right) - x_j \qquad 2 \leq j \leq n-1 \\
g_n(x_{n-1}, \beta, \tau) & = \Phi \left(\tau, x_{n-1}\right) - Q_\partial(\beta)
\end{align*}
for fixed constants $\setof{t_1,\dotsc, t_{n-1}}$.

We want to find a root of $g$ via Newton's method. So, next we compute a formula for $Dg$. First, we let $M(j,k)$ denote the space of $j \times k$ matrices.  We compute the derivative of each component function as block matrix elements in $M(4, 4n-1)$. 

\subsubsection*{First block}
\[
Dg_1(\alpha, x_1) = 
\begin{pmatrix}
\pardiff{}{\alpha} \Phi(t_1, P_\partial(\alpha)) & -\operatorname{I_d} & 0
\end{pmatrix}
\]
where the first block is a column vector since
\[
\pardiff{}{\alpha} \Phi(t_1, P_\partial(\alpha)) = D\Phi(t_1, P_\partial(\alpha)) \cdot DP(e^{i\alpha}, e^{-i \alpha}) \cdot  \left(i e^{i\alpha}, -i e^{-i\alpha}\right) \in M(4,1)
\]
and $I_d \in M(4,4)$ and the last block of zeros is in $M(4, 4n-6)$.

\subsubsection*{Middle blocks}
For $1 < j < n$, we have
\[
Dg_j(x_{j-1}, x_j) = 
\begin{pmatrix}
0 & D \Phi(t_j, x_{j-1}) & -\operatorname{I_d} & 0
\end{pmatrix}
\]
where the first block of zeros is in $M(4, 4j-7)$ and the last block in $M(4, 4n - 4j - 2)$, and the two center blocks are in $M(4,4)$. 

\subsubsection*{Last block}
Finally, we have 
\[
Dg_n(x_{n-1}, \beta, \tau) = 
\begin{pmatrix}
0 & D \Phi(\tau, x_{n-1}) & f \left(\Phi(\tau, x_{n-1})\right) & - \pardiff{}{\beta} Q_\partial(\beta)
\end{pmatrix}
\]
where the first block of zeros is in $M(4, 4n-7)$, the second block in $M(4,4)$, the third block in $M(4,1)$, and the last block is a column vector since
\[
\pardiff{}{\beta} Q_\partial(\beta) = -DQ(e^{i \beta}, e^{-i \beta}) \cdot (i e^{i\beta}, -ie^{-i \beta}) \in M(4,1)
\]

\subsection{The variational equation}
In order to iterate Newton's method on $g$ defined above, we need to compute expressions for the derivative of the flow, $D\Phi(t, x)$. This is obtained by solving the variational equation which requires computing the derivative of the vector field. We do this only for the original CR3BP since the numerical method will not require automatic differentiation. 
We write the CR3BP as a vector field, $f: \rr^4 \to \rr^4$ defined by 
\[
f(x,p,y,q) = 
\begin{pmatrix}
p \\
2q + \left(1 - \mu\right)\left(x - \mu\right) + \mu\left(x + 1 - \mu\right) - \frac{\paren{1 - \mu} \paren{x - \mu}}{d_1^3} - \frac{\mu \paren{x + 1 - \mu}}{d_2^3} \\
q \\
 - 2p + y - \frac{\paren{1-\mu}y}{d_1^3} - \frac{\mu y}{d_2^3} \\
\end{pmatrix}
\]
where 
\[
d_1 = \sqrt{\paren{1 - \mu}^2 + y^2} \qquad d_2 = \sqrt{\paren{x + 1 - \mu}^2 + y^2}
\]

We first compute partials for $d_i$ which are given by 
\begin{align*}
	\pardiff{}{x} d_1 & = \frac{x - \mu}{d_1} \\
	\pardiff{}{y} d_1 & = \frac{y}{d_1} \\
	\pardiff{}{x} d_2 & = \frac{x + 1 - \mu}{d_2} \\
	\pardiff{}{y} d_2 & = \frac{y}{d_2} 
\end{align*}
Now, we have the partials for $f_2, f_4$ given by 
\[
\pardiff{f_2}{x} = 1 - \paren{1 - \mu} \paren{\frac{1}{d_1^3} - \frac{3 (x- \mu)^2}{d_1^5}} - \mu \paren{\frac{1}{d_2^3} - \frac{3 (x + 1 - \mu)^2}{d_2^5}}
\]
\[
\pardiff{f_2}{y} = \pardiff{f_4}{x} = 3 y \paren{\frac{(1- \mu)(x- \mu)}{d_1^5} + \frac{\mu (x + 1 - \mu)}{d_2^5}}
\]
\[
\pardiff{f_4}{y} = 1 - \paren{1 - \mu} \paren{\frac{1}{d_1^3} - \frac{3 y^2}{d_1^5}} - \mu \paren{\frac{1}{d_2^3} - \frac{3y^2}{d_2^5}}
\]
where the equality in the middle pair is due to the fact that the coordinates of $f$ arise from partial derivatives of the potential, $\Omega \in C^2$. So in fact these coordinates are nothing more than the Hessian of $\Omega$ at $(x,y) \in \rr^2$. Specifically, we write
\[
D^2 \Omega(x,y) := H = 
\begin{pmatrix}
\pardiff{f_2}{x}  & \pardiff{f_2}{y}  \\
\pardiff{f_4}{x}  & \pardiff{f_4}{y}  \\
\end{pmatrix}
= 
\begin{pmatrix}
H_1 & H_2 \\
H_2 & H_3
\end{pmatrix}
\]
With these computed, the Jacobian for $f$ reads
\[
Df(x,p,y,q) = 
\begin{pmatrix}
0 & 1 & 0 & 0 \\
H_1 & 0 & H_2  & 2 \\
0 & 0 & 0 & 1 \\
H_2 & -2 & H_3 & 0 
\end{pmatrix}
\]
\end{document}

